% © 2025 Ing. David Jaroš — CC BY-NC-ND 4.0
%
% This work is licensed under a Creative Commons Attribution-NonCommercial-NoDerivatives
% 4.0 International License (CC BY-NC-ND 4.0).

\documentclass[11pt,a4paper]{article}
\usepackage{amsmath, amssymb, amsthm}
\usepackage{hyperref}
\usepackage{geometry}
\geometry{margin=2.5cm}

\newtheorem{theorem}{Theorem}
\newtheorem{proposition}{Proposition}
\newtheorem{definition}{Definition}
\newtheorem{remark}{Remark}

\title{Canonical Field Equations\\
\large Unified Biquaternion Theory — Canonical Theory Series}
\author{Ing.\ David Jaroš}
\date{2025}

\begin{document}
\maketitle

\begin{abstract}
This document collects the canonical field equations of UBT. There are three
levels: (1) the master biquaternionic equation encoding all sectors,
(2) the Einstein field equations recovered as the real-sector limit,
and (3) the matter/gauge field equations. Each equation is cross-referenced
to the primary source file and labelled with its proof status.
\end{abstract}

\section{Master Field Equation}

The master equation of UBT is the biquaternionic wave equation
(``T-shirt formula''):
\begin{equation}\label{eq:master}
  \boxed{
    \nabla^\dagger\nabla\,\Theta(q,\tau)
    \;=\;
    \kappa\,\mathcal{T}(q,\tau)
  }
\end{equation}
where:
\begin{itemize}
  \item $\Theta(q,\tau) \in \mathcal{B} = \mathbb{H}\otimes_\mathbb{R}\mathbb{C}$
        is the fundamental biquaternionic field
  \item $\nabla^\dagger\nabla$ is the biquaternionic Laplacian
        (conjugate gradient squared)
  \item $\kappa = 8\pi G/c^4$ is the gravitational coupling constant
  \item $\mathcal{T}(q,\tau)$ is the biquaternionic source tensor
  \item $q \in \mathbb{R}^{1,3}$ is the spacetime point,
        $\tau = t + i\psi$ is complex time
\end{itemize}

\noindent\textbf{Status}: Axiom (Layer 0/1); source of all UBT field equations.\\
\noindent\textbf{Source}: \texttt{canonical/fields/theta\_field.tex}, \S3.

\section{Einstein Field Equations (Real Sector)}

\begin{theorem}[Einstein Equations from UBT]
In the real sector (constant phase $\varphi$, $\psi \to 0$ limit), the master
equation \eqref{eq:master} projected to the real metric
$g_{\mu\nu} = \mathrm{Re}(\mathcal{G}_{\mu\nu}[\Theta])$ becomes:
\begin{equation}\label{eq:Einstein}
  G_{\mu\nu}
  \;=\;
  R_{\mu\nu} - \tfrac{1}{2}g_{\mu\nu}R
  \;=\;
  \kappa\,T_{\mu\nu},
\end{equation}
where $G_{\mu\nu}$ is the Einstein tensor and $T_{\mu\nu}$ is the
stress-energy tensor derived (not postulated) from the matter sector of
$\mathcal{T}$.
\end{theorem}

\noindent\textbf{Status}: PROVED [L1] via Hilbert variation.\\
\noindent\textbf{Source}: \texttt{canonical/geometry/gr\_as\_limit.tex}, \S3--\S4;
\texttt{consolidation\_project/GR\_closure/GR\_chain\_summary.tex}, Steps 1--5;
\texttt{canonical/bridges/GR\_chain\_bridge.tex}.

\begin{remark}[GR Embedding]
Equation \eqref{eq:Einstein} is identical to Einstein's field equations.
UBT \emph{embeds} GR as the real limit of a richer biquaternionic geometry;
it does not replace GR. All experimental confirmations of GR validate UBT's
real sector.
\end{remark}

\section{Stress-Energy Tensor}

\begin{proposition}[Derived Stress-Energy]
The stress-energy tensor $T_{\mu\nu}$ in \eqref{eq:Einstein} is derived from
the matter Lagrangian by the Hilbert formula:
\begin{equation}
  T_{\mu\nu}
  \;=\;
  -\frac{2}{\sqrt{-g}}\,\frac{\delta(\sqrt{-g}\,\mathcal{L}_{\mathrm{matter}})}{\delta g^{\mu\nu}}.
\end{equation}
It is not postulated — it follows from the variational principle applied
to the matter and gauge sectors of $\mathcal{L}_{\mathrm{UBT}}$.
\end{proposition}

\noindent\textbf{Status}: PROVED [L1].\\
\noindent\textbf{Source}: \texttt{canonical/geometry/stress\_energy.tex};
\texttt{consolidation\_project/GR\_closure/step4\_offshell\_Tmunu.tex}.

\section{Electromagnetic Field Equations (Maxwell)}

In the U(1) gauge sector, varying the action with respect to $A_\nu$ yields:
\begin{align}
  \partial_{[\lambda}F_{\mu\nu]} &= 0 \label{eq:Maxwell1} \\
  \partial_\mu F^{\mu\nu} &= j^\nu \label{eq:Maxwell2}
\end{align}
Equations \eqref{eq:Maxwell1}--\eqref{eq:Maxwell2} are Maxwell's equations.

\noindent\textbf{Status}: PROVED [L1].\\
\noindent\textbf{Source}: \texttt{canonical/interactions/qed.tex}, \S3;
\texttt{canonical/bridges/Maxwell\_limit\_bridge.tex}.

\section{QED Field Equations}

The Dirac equation for the matter field in the U(1) background:
\begin{equation}\label{eq:Dirac}
  (i\gamma^\mu D_\mu - m)\,\psi \;=\; 0,
  \qquad
  D_\mu = \partial_\mu - ieA_\mu,
\end{equation}
emerges from the matter sector of the UBT Lagrangian
$\mathrm{Tr}[(D_\mu\Theta)^\dagger(D^\mu\Theta)]$
in the fermionic limit.

\noindent\textbf{Status}: PRESENT [P] — Dirac equation follows from standard
QFT applied to the U(1) sector; full biquaternionic derivation of the
Dirac spinor structure is partial.\\
\noindent\textbf{Source}: \texttt{canonical/interactions/qed.tex}, \S4.

\section{Yang--Mills Field Equations}

For the non-abelian gauge sectors (SU(3) and SU(2)):
\begin{equation}\label{eq:YM}
  D_\mu F^{a\,\mu\nu} \;=\; j^{a\,\nu},
\end{equation}
where $F^a_{\mu\nu} = \partial_\mu A^a_\nu - \partial_\nu A^a_\mu
+ g_s f^{abc} A^b_\mu A^c_\nu$ and $f^{abc}$ are the structure constants
of SU(3) or SU(2).

\noindent\textbf{Status}: Yang--Mills equations derived from gauge sector
of UBT Lagrangian. SU(3) fully embedded [L0]; full dynamical content
(confinement, asymptotic freedom) is partial [P].\\
\noindent\textbf{Source}: \texttt{canonical/interactions/sm\_gauge.tex};
\texttt{consolidation\_project/appendix\_E\_SM\_QCD\_embedding.tex}.

\section{Off-Shell Closure (Open Problem)}

\begin{remark}[Open Problem: Off-Shell Θ-Only Closure]
Step 6 of the GR chain — deriving the Einstein equations purely from $\Theta$
without any explicit matter source $\mathcal{T}$ — has a structural obstruction
and remains an open problem (Layer 2). This does \emph{not} invalidate Steps 1--5.

\noindent\textbf{Source}: \texttt{canonical/geometry/gr\_completion\_attempt.tex};
\texttt{canonical/bridges/GR\_chain\_bridge.tex}, \S6;
\texttt{AUDITS/copilot\_repo\_verification\_and\_gap\_report.md}.
\end{remark}

\section{Summary Table}

\begin{center}
\begin{tabular}{llll}
\hline
\textbf{Equation} & \textbf{Sector} & \textbf{Source File} & \textbf{Status} \\
\hline
Master eq.\ \eqref{eq:master} & All & \texttt{theta\_field.tex} \S3 & Axiom \\
Einstein eq.\ \eqref{eq:Einstein} & Gravity & \texttt{gr\_as\_limit.tex} & PROVED [L1] \\
$T_{\mu\nu}$ derivation & Gravity & \texttt{stress\_energy.tex} & PROVED [L1] \\
Maxwell eqs.\ \eqref{eq:Maxwell1}--\eqref{eq:Maxwell2} & EM & \texttt{qed.tex} \S3 & PROVED [L1] \\
Dirac eq.\ \eqref{eq:Dirac} & Matter & \texttt{qed.tex} \S4 & PRESENT [P] \\
Yang--Mills eq.\ \eqref{eq:YM} & Gauge & \texttt{sm\_gauge.tex} & PRESENT [P] \\
Off-shell Θ-closure & Gravity & \texttt{gr\_completion\_attempt.tex} & OPEN [L2] \\
\hline
\end{tabular}
\end{center}

\section{Cross-References}

\begin{itemize}
  \item Axioms: \texttt{THEORY/canonical/canonical\_axioms.tex}
  \item Action: \texttt{THEORY/canonical/canonical\_action.tex}
  \item Metric: \texttt{THEORY/canonical/canonical\_metric\_definition.tex}
  \item Projection rules: \texttt{THEORY/canonical/canonical\_projection\_rules.tex}
  \item GR chain: \texttt{canonical/bridges/GR\_chain\_bridge.tex}
  \item QED/Maxwell: \texttt{canonical/bridges/Maxwell\_limit\_bridge.tex}
  \item Gauge sector: \texttt{canonical/bridges/gauge\_emergence\_bridge.tex}
\end{itemize}

\end{document}
